\chapter{绪论}

\section{研究背景}

多重重夸克偶素末态是研究高能强相互作用、重夸克偶素产生机制以及多部分子相互作用的重要窗口。
相较于单个或成对重夸克偶素末态，\(J/\psi + J/\psi + \phi\) 联合产生兼具更复杂的末态拓扑和更低的产生概率，
因此对触发、重建、组合判别和统计建模都提出了更高要求。

从 \texttt{struct-ref/} 中的公开 CMS 论文与分析资料可以看到，相关工作的论文组织通常沿着
“物理动机、实验与数据、选择、效率、拟合、系统误差、结果”的顺序展开。
本文据此构建适合本科综合论文训练的章节结构，并针对 \(J/\psi + J/\psi + \phi\) 末态做针对性展开。

\section{研究问题与目标}

本文关注的核心问题是：
在 CMS 实验数据中，如何建立一套可复查的 \(J/\psi + J/\psi + \phi\) 候选事例选择和结果呈现框架，
并据此组织毕业论文写作。

围绕这一问题，本文的写作目标包括：

\begin{itemize}
  \item 说明该末态与多部分子相互作用、重夸克偶素联合产生研究之间的关系；
  \item 梳理 CMS 实验中与本分析直接相关的探测器、触发、数据样本和模拟样本信息；
  \item 给出 \(J/\psi\) 与 \(\phi\) 候选的重建和选择流程，以及事件级组合策略；
  \item 预留效率修正、信号提取、系统误差和结果汇总的位置，以便后续填入最终分析输出。
\end{itemize}

\section{本文结构}

本文余下各章按 CMS 实验分析的一般流程组织。
第二章介绍物理动机、相关理论背景和可借鉴的公开结果结构；
第三章介绍 CMS 实验装置以及本分析所需的数据与模拟样本；
第四章说明候选粒子的重建和事件选择；
第五章预留效率修正、信号提取与系统误差评估框架；
第六章用于汇总结果、交叉检验和物理讨论；
第七章给出结论与后续展望。
